
The constraint polyhedron is represented in the figure below:

\begin{center}
\begin{tikzpicture}[scale=1.4]
  \begin{axis}[ xlabel={$x_1$}, ylabel={$x_2$}
      ,axis on top=true
      ,axis equal
      ,axis lines=middle
      ,samples=41
      ,xtick = {1}
      ,ytick = {1}
      ,domain=0:1.5
      ,xmin=-1.3,xmax=1.3
      ,ymin=-0.1,ymax=1.2
,legend pos=south east
]
\addplot [thick,color=gray!20,fill=gray!20, 
                    fill opacity=0.05]coordinates {
  ( -1.3,0) 
  ( -1.3,1)
  ( 1.3,1)
  (1.3,0)
};
\addplot[color=black, thin,domain=-1.2:1.2] {0}
node[pos=0.2,sloped,above] {\tiny $x_2=0$};
\addplot[color=black,thin,domain=-1.2:1.2] {1}
node[pos=0.2,sloped,below] {\tiny$x_2=1$};

\node [circle,fill=black,draw,scale=0.2,label=above right:$A$] at (0,1) {A};
\node [circle,fill=black,draw,scale=0.2,label=below right:$B$] at (0,0) {B};

\end{axis}
\end{tikzpicture}
\end{center}

Note that this polyhedron is $\R^2$ has no vertex. In order to write
the problem in standard form, we need to replace $x_1$ by the
difference of two non negative variables
\[
x_1 = x_1^+ - x_1^-,
\]
and we need to introduce a slack variable for the constraint $x_2 \leq 1$.
We obtain
\[
\min_{x \in \R^4} x_1^+ - x_1^-  + x_2
\]
subject to
\begin{align*}
  x_2 + x_3 &= 1, \\
  x_1^+, x_1^-, x_2, x_3 &\geq 0.
\end{align*}
In matrix form, we have
\[
\min_{x \in \R^4} c^T x
\]
subject to
\begin{align*}
  Ax &= b, \\
  x &\geq 0,
\end{align*}
where
\begin{align*}
  A &= \begin{pmatrix}
    0 & 0 & 1 & 1 
  \end{pmatrix}, \\
  b &= 1, \\
  c &= \rvect{1 \\ -1 \\ 1 \\ 0},\\
  x &= \cvect{x_1^+ \\ x_1^- \\ x_2 \\ x_3}.
\end{align*}
This is a problem with $n=4$ variables and $m=1$
constraints. Therefore, there are four potential bases.
\begin{enumerate}
\item $x_1^+$ in the basis. In this case $B=0$. It is not
  invertible. Therefore, it is not a valid basis. 
\item $x_1^-$ in the basis. In this case $B=0$. It is not
  invertible. Therefore, it is not a valid basis.
\item $x_2$ is in the basis. In this case, $B=1$ and
  \[
  x_2=B^{-1}b=1.
  \]
  It is feasible. The vertex
  \[
\cvect{0 \\ 0 \\ 1 \\ 0},
\]
corresponds to point A on the figure. Note that it is a vertex of
the polyhedron in $R^4$, even it does not correspond to a vertex of
the original polyhedron in $\R^2$.
\item $x_3$ is in the basis. In this case, $B=1$ and
  \[
  x_3=B^{-1}b=1.
  \]
  It is feasible. The vertex
  \[
\cvect{0 \\ 0 \\ 0 \\ 1},
\]
corresponds to point B on the figure. Again, it is a vertex of
the polyhedron in $R^4$, even it does not correspond to a vertex of
the original polyhedron in $\R^2$.
\end{enumerate}



